% !PRINTED 166
\weberpage{166}
%<<CONT
die $\mu$ Producte $A_1F_1, A_2F_2, \ldots, A_\mu F_\mu$ alle von gleichem Grade, dem Grade von $F$ sind. Nat\"urlich aber wird im Allgemeinen nicht gefordert, dass umgekehrt alle Functionen von der Form $A_1F_1+A_2F_2+\cdots+A_\mu F_\mu$ bei beliebigen $A_i$ zu dem Systeme $\mathfrak S$ geh\"oren\footnote{Dies findet nur bei besonderen Systemen $\mathfrak S$ statt, die man Moduln nennt.}.

Der Satz, den wir zu beweisen haben, ist evident, wenn es sich um Functionen einer einzigen Ver\"anderlichen $x_1$ handelt. Denn dann sind alle Formen eines Systemes $\mathfrak S$ Potenzen der Variablen $x_1$ mit nicht negativen Exponenten und mit irgend welchen constanten Co\"efficienten multiplicirt. Identisch verschwindende Functionen brauchen wir nicht zu ber\"ucksichtigen. Nehmen wir dann f\"ur $F_1$ eine dieser Functionen von \gesperrt{m\"oglichst niedrigem Grade}, so kann jede andere Function von $\mathfrak S$ in der Form eines Productes $A_1F_1$ dargestellt werden, worin $A_1$ ebenfalls eine Potenz von $x_1$ mit nicht negativem Exponenten und constantem Co\"efficienten ist.

Um also durch Anwendung der vollst\"andigen Induction zum allgemeinen Beweise zu gelangen, nehmen wir zun\"achst an, der Satz I. sei als richtig erwiesen f\"ur jedes System $\mathfrak S_0$ von Formen von $n$ Variablen $x$ und betrachten zun\"achst ein System $\mathfrak S_r$ von Formen $F$, die ausser den $x$ noch eine $(n+1)^{\text{te}}$ Variable $y$, aber nicht in h\"oherer als der $r^{\text{ten}}$ Potenz enthalten, wenn $r$ irgend eine positive ganze Zahl ist. Jede Function $F$ l\"asst sich dann auf eine einzige Art in die Form setzen:
\weq{2}{F = y^r\varphi + \psi,}
worin die Variable $y$ in $\varphi$ gar nicht mehr und in $\psi$ h\"ochstens bis zur $(r-1)^{\text{ten}}$ Potenz vorkommt.

Wenn $F$ das System $\mathfrak S_r$ durchl\"auft, so durchl\"auft $\varphi$ ein gewisses System $\mathfrak S_0$, das sich nach unserer Voraussetzung durch eine Basis darstellen l\"asst, nehmen wir an in der Form:
\weq{3}{\varphi = a_1\varphi_1 + a_2\varphi_2 + \cdots + a_\mu\varphi_\mu.}
Da nun $\varphi_1,\varphi_2,\ldots,\varphi_\mu$ zu dem Systeme $\mathfrak S_0$ geh\"oren, so giebt es Functionen $F_1, F_2,\ldots,F_\mu$ in $\mathfrak S_r$, so dass
\weq{4}{F_1=y^r\varphi_1+\psi_1,\ F_2=y^r\varphi_2+\psi_2,\ \ldots,\ F_\mu=y^r\varphi_\mu+\psi_\mu,}
und dass $y$ in $\psi_1,\psi_2,\ldots,\psi_\mu$ h\"ochstens bis zur $(r-1)^{\text{ten}}$ Potenz vorkommt.
